\documentclass[12pt]{article}
\usepackage[utf8]{inputenc}
\usepackage[T1]{fontenc}
\usepackage{hyperref}
\title{Dharumbal Language Resources for K-12 Curriculum Development}
\author{Compiled by Claude Sonnet 5.5 (claude-sonnet-5-5) for Andr\'as Kornai}
\date{}
\begin{document}
\maketitle

\section{Introduction}
Dharumbal (Darumbal; ISO 639-3 xgm) is a Pama-Nyungan language of coastal central Queensland around Rockhampton. It is classed as dormant. An Endangered Languages Project entry lists it at the ``Awakening'' stage and quotes a community member who says revitalisation is under way. Community work exists, but I found no school curriculum.

\section{Primary Educational Resources}
The State Library of Queensland publishes a Dharumbal body chart with body-part words: \url{https://content.slq.qld.gov.au/sites/default/files/0005-211883-gamburu-gamu-biyu-dharumbal.pdf}.

A federal media release says the Darumbal People Aboriginal Corporation received funding to make story books, flashcards, website material and other learning resources. I did not see those materials themselves: \url{https://minister.infrastructure.gov.au/burke/media-release/nurturing-first-nations-languages-and-arts}. A grant record on the same project is at \url{https://www.grants.gov.au/Ga/Show/c67fd8e8-3bc0-4a78-a36a-e14248b24542}.

\section{Linguistic and Reference Materials}
Angela Terrill's 2002 book \emph{Dharumbal: the language of Rockhampton} is listed as a reference on the Wikipedia page for the Darumbal people: \url{https://en.wikipedia.org/wiki/Darumbal}. The State Library of Queensland holds early vocabularies, including Walter Roth's 1898 word lists. A related State Library of Queensland post is at \url{https://www.slq.qld.gov.au/blog/rockhampton-indigenous-languages-workshop}.

\section{Religious and Community Resources}
The Darumbal People Aboriginal Corporation leads the revitalisation work. I saw no website for it in the results. The Endangered Languages Project entry is at \url{https://endangeredlanguages.com/node/26564}.

\section{Digital and Audio Resources}
I found no audio or video resources specific to Dharumbal.

\section{Curriculum Development Context}
The language is not known to be taught in schools. The corporation's funded materials (story books, flashcards, website) could be the base for a school programme if the community agrees. Early word lists and Terrill's work are the main documentation. Current speaker numbers and the status of the funded materials were not clear from search results.
\end{document}